% Figure 2
\begin{figure*}[!htp]
    \centering
    \includegraphics[width=0.8\linewidth]
    %{figs/Fig2.jpg}
    {figs/Method.jpg}  
    \caption{\textbf{\NAME{}, our approach for self-supervised denoising using spatiotemporal difference and identity reconstruction}. $I_t$, $I_{t-T}$, and $I_{t-2T}$ are video frames at times t (current frame), t-T, and t-2T. These frames are input to an appearance encoder $\Phi$. The resulting feature representations are input to a spatiotemporal bottleneck $\Theta$ that compresses the 3 appearance features into a single spatiotemporal feature representation. Our model then predicts the reconstruction target, defined in \cref{eqn:reconstruction_target} in \cref{sec:reconstruction}, using the reconstruction decoder $\Psi$. The loss, defined in \cref{eqn:recon_loss}, is calculated and backpropagated through all networks. The architecture is discussed in more detail in \cref{sec:method_arch}.}
    \label{fig:fig2}
\end{figure*}

\section{Related Work}
\label{sec:related_work}

Several works in the image and video denoising space require paired clean/noisy imagery for training, evaluation, or both, limiting their applicability in real-world denoising scenarios \cite{zhang2017beyond, li2023ntire, jiang2025efficient, marcosmorales2023evaluating}. We, notably, do not, and we focus this section on existing self-supervised and unsupervised approaches.

\noindent{}\textbf{Self-supervised and unsupervised image (frame-by-frame) denoising.} Several approaches use variants of blind-spot networks or pixel-wise masking to denoise imagery, including frame-by-frame video. Noise2Self \cite{batson2019noise2self} and Noise2Void (N2V) \cite{krull2019noise2void} train on noisy images without requiring clean targets or paired data. N2V trains a blind-spot network to predict masked pixels' intensities from neighboring pixels. Others \cite{laine2019high, batson2019noise2self, jang2023selfsupervised} refrain from masking pixels via structural blind-spot networks with half-plane receptive-field U-Nets \cite{ronneberger2015unet}.  \citet{jang2023selfsupervised} use a conditional blind-spot network and a loss that regularizes denoised images without masking input pixels. Neighbor2Neighbor \cite{huang2021neighbor2neighbor} proposes a self-supervised loss between two sub-sampled images. In general, noise in real-world imagery, including acoustic imagery, has unknown or non-stationary statistics that are spatially correlated, violating assumptions of pixel-wise independence. 

\noindent{}\textbf{Video Denoising with Flow Estimation}
% Video denoising techniques have adopted non-linear filtering \cite{milanfar2012tour}, sparse prior methods \cite{simoncelli1996noise}, and non-local means \cite{lebrun2013nonlocal, arias2018video} from image denoising methods. 
Some video denoising methods leverage videos' spatiotemporal structure by using optical flow for motion compensation \cite{xue2019video, tassano2020dvdnet}. DVDnet \cite{tassano2020dvdnet} uses calculated flow-estimates to manually warp frames, align their contents, and process them collectively with a CNN. VDFlow \cite{wang2020vdflow} jointly learns video deblurring and optical flow. These approaches target video where one component is moving (\eg an object or the camera). However, when objects' motion affects the background's motion, or when objects are small, we find that leveraging optical flow enlarges objects' motion and makes it difficult to distinguish multiple closely-located objects. An example is in \cref{fig:RAFT_on_cfc} in the Supp Mat.


\noindent{}\textbf{Video Restoration and Denoising with Spatial and Temporal Consistency}
UDVD \cite{Sheth_2021_ICCV} uses a patch-wise noise-to-noise training strategy to predict clean frames by estimating masked pixels from adjacent neighborhoods of noisy frames. It relies on temporal redundancy across frames and does \textit{not} rely on clean imagery, making it appealing for in-the-wild data. However, UDVD assumes gaussian noise, while underwaters sonar and ultrasound videos have pink and speckle noise  \cite{zhang2021seasonal_ambient_noise, zhou2023sonar_speckle, despeckling2013}. UDVD's patchwise nature may cause the model to overwhelmingly learn background when most of the video is spatially and temporally dominated by background. UMVD \cite{UMVD} extends UDVD to focus on microscopy data. It uses frame-level augmentations and a reconstruction loss that predicts each frame from its noisy temporal neighbors. It assumes that the signal is consistent across frames while the noise is not, and it assumes smooth object motion. However, like UDVD, it is also patch-wise. 

AverNet \cite{zhao2024avernet}, another self-supervised video restoration model, similarly relies on temporal consistency to fix time-varying unknown degradations. It uses two modules: 1.) prompt-guided alignment to line up video frames at the pixel level and 2.) prompt-conditioned enhancement, which restores each frame by adapting to the image's degredation. The method was tested and performed well on videos with varying levels of gaussian, poisson, and speckle noise, gaussian and resizing blur, and jpeg and video compression. 

LG-BPN \cite{wang2023lgbpn}, another patch-wise denoising method, uses two main ideas: local guidance and internal consistency. Local guidance preserves image structures like edges or textures that are consistent even with noise. Internal consistency assumes nearby patches in space and time share similar structure. However, LG-BPN struggles with long-range temporal dependency and highly-structured noise, such as sensor artifacts. RVRT \cite{liang2022rvrt} is a transformer-based model for denoising, deblurring, and super-resolution. It processes videos frame-by-frame, but models long-term dependencies (remembering what it learned from frames that are further apart) using spatiotemporal attention. It, however, requires very large memory even for inference. 



% \subsection{Self-supervised learning in video}
% Self-supervised learning in video has been broadly organized by four learning objectives: (1) pretext learning, (2) generative learning, (3) contrastive learning, (4) cross-modal agreement \cite{schiappa2023sslvsurvey}. A pretext task is a predefined task for a network to solve to learn a high-level representation that can be used for later downstream tasks. Examples include but are not limited to playback speed, temporal ordering, solving jigsaw problems, and classifying augmentations. Generative learning in this space often involves frame prediction via GANs or MAEs \cite{sun2022bkind, schiappa2023sslvsurvey, zhu2020CVPR, he2022MAE, Feichtenhofer2022maeSpatioTemporal, Srivastava2015unsuplstm, tulyakov2018MoCoGAN, vondrock2016vidgen, Wei2021MaskedFP}. An encoder takes an input frame or set of frames and projects it to a latent space, then a decoder takes the latent features and projects them to the original space with the goal of reconstructing an original data point. This could be generating a masked frame, generating a masked voxel, predicting the importance of each frame, generating motion maps, or generating superresolution frames \cite{schiappa2023sslvsurvey}. We leverage this space in self-supervised learning to denoise our frames. 
% Contrastive learning is a self-supervised way to push positive input pairs near and negative input pairs further away. Common strategies make positive input pairs 2 different frames from the same video, while negative input pairs are 2 frames from different videos \cite{schiappa2023sslvsurvey}. Cross-modal agreement pertains to when there are multiple modalities (\ie video + text, video + audio). Given that we only have video, this does not apply to our use case. 

\noindent{}\textbf{Denoising autoencoders (DAEs)} were originally introduced to learn robust representations. During training, DAEs intentionally add noise to input data and learn to reconstruct the original uncorrupted signal. mDAE \cite{mDAE}, a method for missing data imputation (replacing missing or unavailable data), improves performance on several datasets. 
% \citet{Zaki2024Explainable} leverages a DAE framework as a preprocesing step to improve the quality of SERS spectra for biomarker quantification and discovery. They generated noisy data by duplicating background measurements at random locations. 
DAEs have also been applied to video tasks. CompDAE \cite{compDAE} explicitly models noise from snapshot compressive imaging measurements in low-light conditions to improve edge detection and depth estimation. TADA \cite{Choi2025TADA} uses an adversarial denoising autoencoder to remove EMG noise from EEG time series data. Our work similarly extends the application of DAEs to spatio-temporal grayscale video denoising; we uniquely combine temporal frames to enhance signal quality while simultaneously addressing the increased noise introduced by this process.

% \noindent{}\textbf{Classical image denoising.}
% A large number of spatial filters from image processing techniques have been applied to image denoising \cite{wiener1949extrapolation, tomasi1998bilateral, pitas1990nonlinear, benesty2010study, yang1996structure, takeda2007kernel, fan2019brief} -- they can be broken down to two types: linear and non-linear filters.
% \textit{Linear spatial filters}. Mean and Gaussian filtering \cite{gonzalez2006meanfilter} reduces Gaussian noise, however, they can over-smooth noisy images \cite{al2015latest}. Weiner filtering \cite{benesty2010study} aims to overcome this drawback, but it can overly blur sharp edges. 
% \textit{Non-linear filters.} With non-linear filters, noise can be reduced without first identifying the noisy pixels. Median-filtering \cite{gonzalez2006meanfilter} replaces each pixel with the median value of its neighboring pixels. Median and bilateral filtering \cite{tomasi1998bilateral} preserves edges while smoothing images to reduce the noise, though bilateral filtering is inefficient \cite{fan2019brief}.

\textbf{Sonar and Ultrasound}
Sonar and ultrasound present challenges to the computer vision community. Key characteristics include pink noise, brighter pixel intensities for non-cavity objects compared to the background, and limited distinctiveness from appearance features. \citet{weld2025standardisationconvexultrasounddata} address ultrasound sensor's variability via geometric analysis and augmentation. Unlike many computer vision datasets based on signal from light intensity, ultrasound, sonar,  lidar, and radar rely on wave echoes. The ``camera'' emits waves that reflect off objects and return to the sensor; distance is measured by the echo's return time. Sonar and ultrasound use sound waves (mainly in liquids), while lidar uses laser pulses and is common in air and land environments \cite{Geiger2013VisionMR, caesar2020nuscenes, chang2019argoverse}. We focus on one sonar dataset and two ultrasound datasets described in more detail in \cref{sec:datasets}. 

% \noindent{}\textbf{Detection by tracking} (DbT) methods locate objects in the first frame of a video and then track them to predict future locations \cite{ravi2024sam2, Aladem2019Combined, Yang2017Character, Ngoc2020New}. Point-tracking is similar, in that pixels are first initialized and then tracking models are trained to follow the pixels in a video \cite{zheng2023point, doersch2022tap, doersch2023tapir, vecerik2023robotap, doersch2024bootstap, koppula2024tapvid, karaev23cotracker, karaev24cotracker3}. If the tracking is robust, detection is solved, even in frames where a detector alone fails.  Tracking, therefore, helps fill the gap where detectors struggle. DbT is used in video object segmentation \cite{ravi2024sam2} to propagate masks over time, in long-term object tracking \cite{Kalal2010Tracking} which re-detects objects when needed, and in Visual Simultaneous Localization and Mapping, which relies on tracking features across frames to infer their presence or position of objects or the camera \cite{Chang2023OTE}. In an underwater setting, point-DbT for object discovery or detection may struggle due to the turbidity of water -- discovering erraneous points since the background has motion. Inspired by DbT approaches, our method adds temporal information directly to the spatial dimension to boost the foreground signal.

% \noindent{}\textbf{Tracking by detection}
% (TbD) \cite{Guo2022, Ciaparrone2020DeepLearning} approaches divide the tracking problem into two steps: first, an object detector predicts objects and their locations in every frame; then, a tracker associates detections from one frame to another, creating trajectories and uniquely identifying objects. While there has been notable progress in TbD recently \cite{Ciaparrone2020DeepLearning}, there are also particular drawbacks. TbD relies heavily on detection performance -- false alarms, detection gaps, and missed object groupings may lead to incomplete tracks. In addition, some recent progress in TbD can be attributed to matching targets with spatially re-identifiable features. However, this is unrealiable in sonar. \citet{cfc2022eccv}'s downstream tracking is a tracking-by-detection approach; in order to improve it, then, we design our method to maximize downstream detection performance by lowering false positive and false negative rates concurrently.